\documentclass[12pt]{article}%
\usepackage{graphicx, verbatim}
\usepackage{amsmath}
\usepackage{amsfonts}
\usepackage{amssymb}
\usepackage{latexsym}
\usepackage{enumerate}
\usepackage[spanish, provide=*]{babel}
\usepackage{palatino}
\usepackage[utf8]{inputenc}


\setlength{\headsep}{-2cm} \setlength{\oddsidemargin}{-0.9cm}
\setlength{\evensidemargin}{1.5cm} \setlength{\textheight}{25cm}
\setlength{\textwidth}{17.7cm} \setlength{\parindent}{0cm}
\parindent=0mm

\pagestyle{empty}

\def\R{\mathbb{R}}
\def\Z{\mathbb{Z}}
\def\C{\mathbb{C}}
\def\N{\mathbb{N}}
\def\A{\mathcal{Z}}
\def\V{\mathcal V}
\def\W{\mathcal W}
\def\F{\mathcal F}
\def\a{\alpha}
\def\vf{\varphi}
\def\re{\operatorname{Re}}
\def\log{\operatorname{Log}}
\def\adj{\operatorname{adj}}
\def\img{\operatorname{Im}}
\def\nu{\operatorname{Nu}}
\newtheorem{exercise}{Ejercicio}
\newcommand{\bej}{\begin{exercise}\rm}
\newcommand{\fej}{\end{exercise}}

% \renewcommand{\labelenumi}{\bf{(\numb{enumi})}}
\renewcommand{\labelenumii}{(\alph{enumii})}

\begin{document}

\pagestyle{empty}%

\label{t1}


\bigskip

%\hfill\textsc{\textbf{Tema }}

\textsc{Nombre y apellido:}

\vskip0.2cm

\textsc{DNI y Libreta Universitaria:}

\vskip0.2cm

\textsc{Carrera:}

\vskip0.2cm

\textsc{Mail:}

\vskip0.4cm


\def\espac{\LARGE$\phantom{\displaystyle\sum\, \,}$}

\begin{tabular}[c]{|c | c | c | c| c|}
\hline
1 & 2 & 3 & 4 & Calif. \\
\hline
\espac& \espac & \espac & \espac & \espac \\
\hline
\end{tabular}

\bigskip
\medskip

\noindent

\hrulefill\ \vskip8pt

\begin{large}\centerline{{\textbf{C\'alculo Num\'erico / Elementos de C\'alculo Num\'erico}}}\end{large}

\vskip.2cm

\begin{large}\centerline{{\textbf{Segundo Parcial - 24 de junio de 2026}}}\end{large}

\hrulefill

\vskip 0.5cm

\begin{exercise}[Matriz de diferenciaci\'on de Chebyshev para ecuaci\'on diferencial ordinaria]
Considere el problema de contorno
\[
-u''(x)+u(x)=1,\qquad x\in[-1,1],\qquad u(-1)=u(1)=0.
\]
Tome los nodos de Chebyshev--Gauss--Lobatto para $N=2$:
\[
x_0=1,\;x_1=0,\;x_2=-1,
\]
y la matriz de diferenciación de Chebyshev de primer orden
\[
D=
\begin{pmatrix}
\tfrac32 & -2 & \tfrac12\\
\tfrac12 & 0 & -\tfrac12\\
-\tfrac12 & 2 & -\tfrac32
\end{pmatrix}.
\]
\begin{enumerate}[(a)]
\item Si $U=(U_0,U_1,U_2)$ es el vector de aproximaciones de la función $u$ en los nodos $x_0, x_1, x_2$, ¿cómo puede aproximarse la derivada segunda de $u$ en los nodos usando $D$ y $U$?
\item Usando la ecuación diferencial, las condiciones de borde $u(-1)=u(1)=0$, y la matriz $D$, obtenga una ecuación para aproximar $u(x_1)$. Resuelva esa ecuación y obtenga la aproximación numérica de $u(x_1)$.
\end{enumerate}
\end{exercise}


\begin{exercise}[Diferencias finitas para reacci\'on-difusi\'on]
Considere
\[
u_t=u_{xx}-u,\qquad 0<x<1,\ t>0,
\]
con condiciones de borde homog\'eneas $u(0,t)=u(1,t)=0$.
Use una malla uniforme $x_j=jh$, $j=0,\dots,M$, y $t^n=n\Delta t$.

Se propone el esquema impl\'icito (Euler hacia atr\'as en tiempo y diferencia
centrada en espacio).
\begin{enumerate}[(a)]
\item Obtenga el sistema lineal tridiagonal $A\,\mathbf u^{n+1}=\mathbf u^n$ que hay que resolver en cada paso de tiempo.
\item Para el problema peri\'odico asociado (es decir, con condiciones de borde $u(0,t)=u(1,t)$), haga un an\'alisis de Fourier y
obtenga el factor de amplificaci\'on $G(\theta)$. ¿El esquema es estable para todo $\Delta t,h>0$?
\end{enumerate}
\end{exercise}


\begin{exercise}[M\'etodo IMEX para ODE]
Considere el problema escalar
\[
y'=-\alpha y+g(y),\qquad y(0)=y_0,
\]
donde $\alpha>0$ y $g:\R\to\R$ es Lipschitz con constante $L_g$. Considere el siguiente esquema IMEX de orden 1 (parte lineal impl\'icita y parte no lineal expl\'icita):
\[
\frac{y^{n+1}-y^n}{\Delta t}=-\alpha y^{n+1}+g(y^n),\qquad n\ge 0.
\]

\begin{enumerate}[(a)]
\item Despeje $y^{n+1}$ y escriba el operador de avance $\Phi_{\Delta t}$ que satisface
$y^{n+1}=\Phi_{\Delta t}(y^n)$.
\item Halle una constante de Lipschitz de $\Phi_{\Delta t}$ en funci\'on de
$\Delta t$, $\alpha$ y $L_g$. Muestre que se puede tomar una cota
$\operatorname{Lip}(\Phi_{\Delta t})\le 1+\Delta t\,L_g,$ y concluya que existe una constante de Lipschitz del m\'etodo independiente de $\alpha$.
\item Para dos soluciones num\'ericas $\{y^n\}$ y $\{z^n\}$ con datos iniciales
distintos, obtenga una cota de estabilidad para $|y^n-z^n|$ hasta tiempo
$T=n\Delta t$.
\end{enumerate}

\end{exercise}


\begin{exercise}[M\'etodo de Newton]
En un circuito con una fuente continua $V_s>0$, una resistencia $R>0$ y un
diodo ideal de Shockley con par\'ametros $I_s>0$ y $V_T>0$, la tensi\'on en el
diodo $v$ satisface
\[
I_s\left(e^{v/V_T}-1\right)=\frac{V_s-v}{R}.
\]
Defina
\[
f(v)=I_s\left(e^{v/V_T}-1\right)-\frac{V_s-v}{R}.
\]
\begin{enumerate}[(a)]
\item Escriba la iteraci\'on de Newton para resolver $f(v)=0$, es decir,
una f\'ormula expl\'icita para $v_{k+1}$ en funci\'on de $v_k$.
\item Verifique que $f$ es estrictamente creciente y convexa en $\R$, y
concluya que la ecuaci\'on tiene una \'unica ra\'iz $v^*$ (de hecho,
$0<v^*<V_s$).
\item Considere $v_0=V_s$. Usando convexidad, muestre que la sucesi\'on de
Newton queda bien definida, cumple $v^*\le v_{k+1}\le v_k$, y por lo tanto
converge a $v^*$.
\end{enumerate}
\end{exercise}




\vskip 2.5cm
\begin{center}
    \textit{En todos los casos debe justificar sus respuestas.}
\end{center}


\end{document}
